\documentclass[10pt]{article}
\usepackage[T1]{fontenc}
\usepackage[utf8]{inputenc}
\usepackage[french]{babel}
\usepackage{amsmath,amssymb,mathrsfs,mathpazo}
\usepackage{geometry}
\geometry{a4paper,margin=2cm}
\usepackage{graphicx,array,tabularx,tabularray,colortbl,multirow,multicol,booktabs,enumitem,pifont,ulem,fancybox,diagbox,cancel}
\usepackage{pstricks,pstricks-add,pst-plot,pst-tree,pst-node,pst-text,pst-func,pst-math,pst-bspline,pst-3dplot,pst-eucl,pst-solides3d}
\let\pscalculate\undefined
\usepackage{pst-bezier}
\usepackage{tikz,pgfplots,tkz-tab}
\usetikzlibrary{arrows,arrows.meta,patterns,calc,positioning,trees,decorations.markings,shapes.geometric,shapes.misc,shapes.arrows,intersections,decorations.pathmorphing,decorations.pathreplacing,backgrounds,shadows,babel}
\pgfplotsset{compat=1.18}
\usepgfplotslibrary{fillbetween}
\usepackage[np]{numprint}
\usepackage{listings,makecell,pas-tableur,twoopt}
\let\barre\undefined
\usepackage[pstricks]{tablvar}
\definecolor{darkgreen}{RGB}{0,100,0}
\providecommand{\red}{\color{red}}
\providecommand{\blue}{\color{blue}}
\providecommand{\green}{\color{green}}
\providecommand{\black}{\color{black}}
\providecommand{\white}{\color{white}}
\providecommand{\gray}{\color{gray}}
\providecommand{\cyan}{\color{cyan}}
\providecommand{\magenta}{\color{magenta}}
\providecommand{\yellow}{\color{yellow}}
\providecommand{\violet}{\color{violet}}
\providecommand{\orange}{\color{orange}}
\providecommand{\euro}{\texteuro}
\providecommand{\EUR}{\texteuro}
\providecommand{\ieme}{\textsuperscript{e}}
\providecommand{\degres}{\ensuremath{^\circ}}
\providecommand{\degree}{\ensuremath{^\circ}}
\providecommand{\vv}[1]{\overrightarrow{#1}}
\providecommand{\Coord}[1]{#1}
\providecommand{\decosix}{}
\providecommand{\starredbullet}{\ensuremath{\bullet}}
\pagestyle{empty}
\setlength{\parindent}{0pt}
\usepackage{adjustbox}\DeclareUnicodeCharacter{2236}{:}\DeclareUnicodeCharacter{2190}{\ensuremath{\leftarrow}}
\providecommand{\R}{\mathbb{R}}
\providecommand{\N}{\mathbb{N}}
\providecommand{\Z}{\mathbb{Z}}
\providecommand{\Q}{\mathbb{Q}}
\providecommand{\D}{\mathbb{D}}
\providecommand{\C}{\mathbb{C}}
\providecommand{\e}{{\mathrm{e}}}
\providecommand{\vect}[1]{\overrightarrow{#1}}
\providecommand{\vectt}[1]{\overrightarrow{\mathrm{#1}}}
\providecommand{\vvt}[1]{\overrightarrow{\mathrm{#1}}}
\providecommand{\np}[2][]{{#2~#1}}
\providecommand{\vv}[1]{{\overrightarrow{#1}}}
\renewcommand{\d}{{\mathrm{d}}}
\providecommand{\degres}{{^\circ}}
\providecommand{\degre}{{^\circ}}
\providecommand{\blue}{\color{blue}}
\providecommand{\red}{\color{red}}
\providecommand{\enskip}{\quad}
\providecommand{\llbracket}{\lbrack\!\lbrack}
\providecommand{\rrbracket}{\rbrack\!\rbrack}
\providecommand{\euro}{{\text{€}}}
\providecommand{\no}{{\text{n°}}}
\providecommand{\interval}[2]{{[#1;#2]}}
\providecommand{\intervFO}[2]{{[#1;#2[}}
\providecommand{\dotfill}{\ldots}

\providecommand{\ds}{\displaystyle}
\providecommand{\textding}[1]{#1}
\providecommand{\Oij}{\(\mathrm{O};\vec{\imath},\vec{\jmath}\)}
\providecommand{\Oijk}{\(\mathrm{O};\vec{\imath},\vec{\jmath},\vec{k}\)}
\providecommand{\Ouv}{\(\mathrm{O};\vec{u},\vec{v}\)}
\providecommand{\intervFF}[2]{\ensuremath{[#1;#2]}}
\renewcommand{\euro}{\ifmmode\text{\texteuro}\else\texteuro\fi}
\DeclareUnicodeCharacter{00B0}{\ifmmode^{\circ}\else\textdegree\fi}
\DeclareUnicodeCharacter{20AC}{\ifmmode\text{\texteuro}\else\texteuro\fi}
% Énoncés : Éducation nationale. Transcription et corrigés : APMEP, Denis Vergès et contributeurs. Découpage et classement : CoopMaths, Sébastien Lozano, Éric Elter et contributeurs. Adaptation de la présentation : Maths Explorer.
% CC BY-SA : https://coopmaths.fr/www/about/
% https://coopmaths.fr/alea/static/bac/2022/tex/bac_2022_09_sujet1_ameriquesud_4.tex
\begin{document}
\noindent\textit{APMEP / CoopMaths. CC BY-SA. Présentation adaptée par Maths Explorer.}\par\medskip
\medskip

Dans la figure ci-dessous, ABCDEFGH est un parallélépipède rectangle tel que 

AB $= 5$, AD $= 3$ et AE $= 2$.

L'espace est muni d'un repère orthonormé d'origine A dans lequel les points B, D et E ont respectivement pour coordonnées (5~;~0~;~0), (0~;~3~;~0) et (0~;~0~;~2).

\begin{center}
\psset{unit=1cm}
\begin{pspicture}(9,5)

\psframe(0.2,0.2)(7.2,3)
\psline(7.2,0.2)(8.4,1.6)(8.4,4.4)(7.2,3)
\psline(8.4,4.4)(1.6,4.4)(0.2,3)
\psline[linestyle=dashed](0.2,0.2)(1.6,1.6)(8.4,1.6)
\psline[linestyle=dashed](1.6,1.6)(1.6,4.4)
\pspolygon[linestyle=dotted,linewidth=1.5pt](0.2,0.2)(8.4,1.6)(3,4.4)
\uput[dl](0.2,0.2){A} \uput[dr](7,0.2){B} \uput[r](8.4,1.6){C}
\uput[ur](1.6,1.6){D} \uput[l](0.2,3){E} \uput[ul](7.2,3){F}
\uput[ur](8.4,4.4){G} \uput[ul](1.6,4.4){H} \uput[u](3,4.4){M}
\end{pspicture}
\end{center}

\begin{enumerate}
\item 
	\begin{enumerate}
		\item Donner, dans le repère considéré, les coordonnées des points H et G.
		\item Donner une représentation paramétrique de la droite (GH).
	\end{enumerate}	
\item Soit M un point du segment [GH] tel que $\vect{\text{HM}} =k\vect{\text{HG}}$ avec $k$ un nombre réel de l'intervalle [0~;~1].
	\begin{enumerate}
		\item Justifier que les coordonnées de M sont $(5k~;~3~;~2)$.
		\item En déduire que $\vect{\text{AM}} \cdot \vect{\text{CM}} = 25k^2  - 25k + 4$.
		\item Déterminer les valeurs de $k$ pour lesquelles AMC est un triangle rectangle en
M.
	\end{enumerate}
\end{enumerate}	

Dans toute la suite de l'exercice, on considère que le point M a pour coordonnées (1~;~3~;~2).

On admet que le triangle AMC est rectangle en M .

On rappelle que le volume d'un tétraèdre est donné par la formule $\dfrac13 \times\text{Aire de la base}  \times h$ où $h$ est la hauteur relative à la base.

\begin{enumerate}[resume]
\item On considère le point K de coordonnées (1~;~3~;~0).
	\begin{enumerate}
		\item Déterminer une équation cartésienne du plan (ACD).
		\item Justifier que le point K est le projeté orthogonal du point M sur le plan (ACD).
		\item En déduire le volume du tétraèdre MACD.
	\end{enumerate}	
\item On note P le projeté orthogonal du point D sur le plan (AMC).

Calculer la distance DP ; en donner une valeur arrondie à $10^{-1}$.
\end{enumerate}

\end{document}
